\section{Problem Definition}
(todo: this section is not good). i want you to focus on memory augmented LLM. where you have LLM, which take external memory to answer question. very formally and mathmatically. 
\paragraph{Conversational Agent.}

We consider an LLM-based agent augmented with an external memory system that persists across interactions:
\begin{equation}
\mathcal{A} = (\pi_\theta,\; \mathcal{X},\; \mathcal{Y},\; \mathcal{M},\; \mathcal{E}),
\end{equation}
\begin{itemize}[nosep]
    \item \textbf{policy} $\pi_\theta$: maps an observation and a memory state to a distribution over action trajectories;
    \item \textbf{input (observation) space}  $\mathcal{X}$: the set of all possible user queries and dialogue the agent may receive;
    \item \textbf{output (response) space} $\mathcal{Y}$ : the set of all possible natural-language responses the agent may produce;
    \item  \textbf{memory space} $\mathcal{M}$: the set of all possible external memory states the agent can read from and write to;
    \item \textbf{external environment interface} $\mathcal{E}$: the set of tools through which the agent can perform actions beyond pure language generation (e.g., a memory retrieval server, a search engine, or a code executor).
\end{itemize}


An \emph{interaction} is a temporally ordered sequence of turns between a user and the agent:
\begin{equation}
\mathcal{I} = (t_1, t_2, \ldots, t_T),
\end{equation}
where $T$ is the total number of turns in the interaction.
At turn $t$ ($1 \le t \le T$), the user provides an input $x_t \in \mathcal{X}$ (a query, instruction, or follow-up). The agent observes $x_t$ together with the current memory state $M_t \in \mathcal{M}$. T

The agent generates an action trajectory as a sequence of \emph{language actions} (reasoning steps, final answers) and \emph{environment actions} (memory queries, tool calls).
The agent emits a final response $y_t \in \mathcal{Y}$ extracted from $\tau_t$. The memory state may be updated: $M_{t+1} = f_{\text{update}}(M_t, \tau_t)$, where $f_{\text{update}}$ encodes any write operations performed during the turn.
The objective is to learn a policy $\pi_\theta$ that controls memory access when and how memory should be queried to improve downstream task performance. 

\paragraph{Memory Structures}

Recent work, notably Search-R1~\citep{jin2025search}  and Memory-R1~\citep{yan2025memory}\footnote{Memory-R1 does not provide a public codebase. We built a publicly available implementation at \url{https://github.com/YifanLiu2/Struct_Memory_R1.git}.}, has shown that RL can effectively train agents to manage flat external memory banks.
A \emph{flat memory state} is an unordered collection of $N$ memory entries:
\begin{equation}
M_t^{\text{flat}} = \{m_1, m_2, \ldots, m_N\}, \quad m_i = (\text{key}_i,\; \text{attributes}_i),
\end{equation}